\vspace{-1.5ex}
\section{Method}
\vspace{-0.5ex}
\label{sec:method}

We propose \name,
a unified VLP framework to learn from noisy image-text pairs.
% with two novel components:
% (1) Mixture of Encoder-Decoder for effective multi-task pre-training and flexible transfer learning;
% (2) Bootstrap Captions for improving the quality of the pre-training corpus. 
This section first introduces our new model architecture MED and its pre-training objectives,
and then delineates CapFilt for dataset bootstrapping. 

\vspace{-1ex}
\subsection{Model Architecture}
\label{sec:med}
\vspace{-0.5ex}

We employ a visual transformer~\cite{vit} as our image encoder,
which divides an input image into patches and encodes them as a sequence of embeddings, with an additional \texttt{[CLS]} token to represent the global image feature.
Compared to using pre-trained object detectors for visual feature extraction~\cite{uniter},
using a ViT is more computation-friendly and has been adopted by the more recent methods~\cite{ALBEF,ViLT}.

In order to pre-train a unified model with both understanding and generation capabilities,
we propose multimodal mixture of encoder-decoder (MED),
a multi-task model which can operate in one of the three functionalities:

\vspace{-0.8ex}
(1) \textbf{Unimodal encoder}, which separately encodes image and text. The text encoder is the same as BERT~\cite{bert}, where a \texttt{[CLS]} token is appended to the beginning of the text input to summarize the sentence.

\vspace{-0.8ex}
(2) \textbf{Image-grounded text encoder}, which injects visual information by inserting one additional cross-attention (CA) layer between the self-attention (SA) layer and the feed forward network (FFN) for each transformer block of the text encoder. A task-specific \texttt{[Encode]} token is appended to the text, and the output embedding of \texttt{[Encode]} is used as the multimodal representation of the image-text pair. 

\vspace{-0.8ex}
(3) \textbf{Image-grounded text decoder}, which replaces the bi-directional self-attention layers in the image-grounded text encoder with causal self-attention layers. A \texttt{[Decode]} token is used to signal the beginning of a sequence, and an end-of-sequence token is used to signal its end.


\vspace{-1ex}
\subsection{Pre-training Objectives}
\vspace{-0.5ex}
We jointly optimize three objectives during pre-training,
with two understanding-based objectives and one generation-based objective.
Each image-text pair only requires one forward pass through the computational-heavier visual transformer,
and three forward passes through the text transformer,
where different functionalities are activated to compute the three losses as delineated below.

\noindent\textbf{Image-Text Contrastive Loss} (ITC) activates the unimodal encoder. It aims to align the feature space of the visual transformer and the text transformer by encouraging positive image-text pairs to have similar representations in contrast to the negative pairs. 
It has been shown to be an effective objective for improving vision and language understanding~\cite{clip, ALBEF}.
We follow the ITC loss by~\citet{ALBEF},
where a momentum encoder is introduced to produce features,
and soft labels are created from the momentum encoder as training targets to account for the potential positives in the negative pairs. 

\noindent\textbf{Image-Text Matching Loss} (ITM) activates the image-grounded text encoder.
It aims to learn image-text multimodal representation that captures the fine-grained alignment between vision and language. 
ITM is a binary classification task,
where the model uses an ITM head (a linear layer) to predict whether an image-text pair is positive (matched) or negative (unmatched) given their multimodal feature.
In order to find more informative negatives,
we adopt the hard negative mining strategy by~\citet{ALBEF},
where negatives pairs with higher contrastive similarity in a batch are more likely to be selected to compute the loss.

\noindent\textbf{Language Modeling Loss} (LM) activates the image-grounded text decoder,
which aims to generate textual descriptions given an image.
It optimizes a cross entropy loss which trains the model to maximize the likelihood of the text in an autoregressive manner.
We apply a label smoothing of 0.1 when computing the loss.
Compared to the MLM loss that has been widely-used for VLP,
LM enables the model with the generalization capability to convert visual information into coherent captions.


In order to perform efficient pre-training while leveraging multi-task learning,
the text encoder and text decoder share all parameters except for the SA layers.
The reason is that the differences between the encoding and decoding tasks are best captured by the SA layers. In particular, the encoder employs \textit{bi-directional} self-attention to build representations for the \textit{current} input tokens, while the decoder employs \textit{causal} self-attention to predict \textit{next} tokens.
On the other hand, the embedding layers, CA layers and FFN function similarly between encoding and decoding tasks,
therefore sharing these layers can improve training efficiency while benefiting from multi-task learning,


% which is qualitatively shown in Figure~\ref{fig:heatmap}.
% For each text token $t_i$,
% the encoder's SA layer builds a representation for $t_i$,
% which is then used by the following CA layer to query the corresponding image features.
% Different from the encoder, 
% the decoder's SA layer builds a representation for the \textit{next} possible token $t_{i+1}$,
% which is used by the CA layer to query the image in search for visual information that corresponds to $t_{i+1}$.
% In addition, encoder uses a bi-directional self-attention mechanism whereas the decoder uses causal self-attention.

\input{figure/CapFilt}

\vspace{-0.5ex}
\subsection{CapFilt}
\vspace{-0.5ex}
Due to the prohibitive annotation cost,
there exist a limited number of high-quality human-annotated image-text pairs $\{(I_h,T_h)\}$ (e.g., COCO~\cite{coco}).
Recent work~\cite{ALBEF,simvlm} utilizes a much larger number of image and alt-text pairs $\{(I_w,T_w)\}$ that are automatically collected from the web.
However, the alt-texts often do not accurately describe the visual content of the images, making them a noisy signal that is suboptimal for learning vision-language alignment. 

We propose Captioning and Filtering (CapFilt),
a new method to improve the quality of the text corpus.
Figure~\ref{fig:capfilt} gives an illustration of CapFilt. 
It introduces two modules: a \textit{captioner} to generate captions given web images,
and a \textit{filter} to remove noisy image-text pairs.
Both the captioner and the filter are initialized from the same pre-trained MED model,
and finetuned individually on the COCO dataset.
The finetuning is a lightweight procedure.

Specifically,
the \emph{captioner} is an image-grounded text decoder.
It is finetuned with the LM objective to decode texts given images. 
Given the web images $I_w$, the captioner generates synthetic captions $T_s$ with one caption per image.
The \emph{filter} is an image-grounded text encoder.
It is finetuned with the ITC and ITM objectives to learn whether a text matches an image.
The filter removes noisy texts in both the original web texts $T_w$ and the synthetic texts $T_s$,
where a text is considered to be noisy if the ITM head predicts it as unmatched to the image.
Finally,
we combine the filtered image-text pairs with the human-annotated pairs to form a new dataset,
which we use to pre-train a new model.












